\section{Related Work}
\label{sec:related}

\para{The same request, different behavior by perceived user.} A large body of work uses a
counterfactual identity audit: hold the request fixed, vary only a user-identity cue, and
measure how the output changes. \citet{vijjini2024safety} define a ``personalization bias''
in which safety guarding and utility shift purely with stated demographics, and find adding
a minor identity typically raises safety while a non-binary identity lowers it --- the
closest analogue to differential protective advice. \citet{gupta2023bias} show that
assigning a socio-demographic persona makes models abstain or err based on identity, with up
to $70\%$ relative accuracy drops and abstention-style refusals concentrated on
disadvantaged personas. \citet{vli2024guardrail} find refusal rates on a fixed sensitive
request swing with a user's politics, age, gender, and race. \citet{childsafety2025} report
$40$--$60$ percentage-point higher safety-defect rates for child versus adult personas on
identical harm topics --- the single most go-to-sleep-relevant result, since it shows models
demonstrably treat a perceived child differently. \citet{position2025} show identity placed
in the system prompt produces stronger, more opaque conditioning than the same cue in the
user turn. Unlike these studies, which vary the person while holding the situation fixed, we
cross a full factorial of situations and personas so we can measure their \emph{relative}
contributions.

\para{Wellbeing and emotional advice conditioned on the user.} Closest to our domain,
intersectional empathy audits hold an emotional situation fixed and vary persona, finding
that empathy and solution-giving shift systematically and stereotypically, with a
topic-to-attribute variance ratio above one for several groups --- the response is driven
more by who the user is than by what they are going
through~\citep{empathy2025}. \citet{memory2025} show that even stored user profiles shift
answers on validated, user-independent emotional-intelligence tests.
\citet{goalpref2026} add a crucial nuance: with no instruction to use the persona, models
still calibrate validation and hedging to the inferred user, but the \emph{direction} is
governed by the situational \emph{role} (more challenging as an advisor, more deferential as
a peer). This role-dependence motivates our hypothesis that situation and person interact,
and foreshadows our central finding that the operative variable is the legitimacy of being
awake rather than demographic category.

\para{Foundations: steerability, sycophancy, persona framing.} \citet{santurkar2023opinions}
establish that persona injection steers model opinions toward a named group, but only
modestly and incompletely. \citet{sharma2023sycophancy} document a different axis of
person-conditioning --- matching the user's stated beliefs --- that our probe must guard
against. Surveys of persona-steered generation formalize the distinction we adopt between
\emph{role-playing} (situation-driven) and \emph{personalization} (person-driven).

\para{The counterweight: don't overclaim.} \citet{eloundou2024fairness}, the largest and
most rigorous such audit, vary only the user's name and find no statistically significant
difference in response quality, with harmful stereotyping rare and $3$--$12\times$ smaller
after RLHF. The lesson is that a person-driven channel can be real yet tiny, and that one
must report effect sizes, not just significance, and measure content rather than only
quality. \citet{differentcues2026} further caution that measured conditioning depends
heavily on which cue and framing are used, motivating our multi-factor design. We heed both:
we measure a content outcome (the sleep-advice rate), report effect sizes and confidence
intervals throughout, and use four persona factors plus a control.

\para{Domain anchor.} Personal-health LLMs are designed to give sleep advice conditioned on
user data~\citep{phllm2024}, but they do not audit identity fairness. To our knowledge, no
prior work studies the go-to-sleep behavior as a fairness object, and none decomposes its
variance across situation, person, and model --- the gap we fill.
